\chapter{Jeux joués}\label{ch:iv2:played-games}

\begin{ivtrains}
Coter des prix à voix haute pendant que de nouvelles informations arrivent. Chaque jeu a ici ses règles et une feuille de score, et un script dans sa réponse. Recopier la feuille de score, ouvrir le script et le couvrir d'une feuille de papier. Ne faire glisser la feuille que jusqu'à la prochaine ligne d'arrêt, écrire sa réponse, puis poursuivre la lecture. Mieux encore, demander à un ami de lire le script à voix haute en tant que recruteur, sans donner le moindre indice.
\end{ivtrains}

\begin{question}[$\star\star\star\star$ \iqroles{QT, Trader, QR} \iqfirm{teneur de marché, hedge fund, crypto} \iqlevel{junior, confirmé}]\label{iq:iv2:played-games:1}
\textbf{Trois cartes et un joueur mieux informé.} Dix cartes numérotées de 1 à 10 sont mélangées, et trois sont distribuées face cachée~: $A$, $B$ et $C$. Vous faites un marché sur leur somme. Un lot acheté à un prix $x$ se règle à la somme moins $x$, et un lot vendu à $x$ se règle à $x$ moins la somme.

Vous cotez quatre fois, à chaque fois sous forme d'un cours acheteur et d'un cours vendeur, par pas de 0.5. Votre cours vendeur doit être exactement 2 au-dessus de votre cours acheteur lors des manches 1 à 3, et exactement 1 au-dessus lors de la manche 4. Aucune carte n'est face visible lors des manches 1 et 2. La carte $A$ est retournée avant la manche 3, et la carte $C$ avant la manche 4.

Votre seule contrepartie est un autre joueur qui a regardé la carte $B$ et qui calcule parfaitement. À chaque cotation, il calcule la somme espérée, sachant $B$ et les cartes visibles. Il vous vend un lot à votre cours acheteur si celui-ci est supérieur à cette valeur, il vous achète un lot à votre cours vendeur si celui-ci est inférieur à cette valeur, et sinon il passe. Vous connaissez cette règle. Prenez environ 20 secondes pour chaque cotation, et essayez de terminer avec le plus d'argent possible.

\medskip
\begin{tabular}{@{}l|l|l|l|l|l@{}}
Manche & Votre valeur juste & Cours acheteur & Cours vendeur & Transaction & Votre position \\ \hline
\ivscoreline{1}
\ivscoreline{2}
\ivscoreline{3}
\ivscoreline{4}
\end{tabular}
\end{question}
\vfill\clearpage

\begin{question}[$\star\star\star\star$ \iqroles{QT, Trader} \iqfirm{teneur de marché, hedge fund} \iqlevel{junior, confirmé}]\label{iq:iv2:played-games:2}
\textbf{Options sur deux dés.} Deux dés équilibrés sont lancés l'un après l'autre, et $U$ est leur somme. Trois instruments se règlent sur $U$~: un contrat à terme versant $U$, un call versant $\max(U - 8, 0)$ et un put versant $\max(8 - U, 0)$. Vous êtes dealer. Vous négociez des lots entiers, vous pouvez détenir n'importe quelle position, et vous êtes évalué sur votre profit final et sur le risque que vous avez pris pour l'obtenir.

Le jeu comporte cinq manches. Lors de la manche 1, vous cotez le call, avec un cours acheteur et un cours vendeur écartés d'au plus 0.30, avant que le moindre dé ne soit lancé, et un client peut traiter avec vous. Lors de la manche 2, vous pouvez échanger des contrats à terme avec un autre dealer. Le dé 1 est lancé avant la manche 3, où vous pouvez de nouveau échanger des contrats à terme. Lors de la manche 4, l'autre dealer affiche des prix pour les trois instruments. Le dé 2 est alors lancé, et lors de la manche 5 tout se règle. Prenez environ 30 secondes pour chaque manche.

\medskip
\begin{tabular}{@{}l|p{7.5cm}|l@{}}
Manche & Ce que vous avez fait & Cash généré \\ \hline
1 & & \\[8pt]\hline
2 & & \\[8pt]\hline
3 & & \\[8pt]\hline
4 & & \\[8pt]\hline
5 & Règlement & \\[8pt]\hline
\end{tabular}
\end{question}
\vfill\clearpage

\begin{question}[$\star\star\star\star$ \iqroles{QT, Trader, QR} \iqfirm{teneur de marché, hedge fund, crypto} \iqlevel{junior, confirmé}]\label{iq:iv2:played-games:3}
\textbf{Enchérir sur un lot que vous ne connaissez qu'à moitié.} Deux dés à dix faces, numérotées de 1 à 10, sont lancés pour chaque lot. Vous voyez le dé $A$, votre rival voit le dé $B$, et aucun des deux ne voit celui de l'autre. Le lot vaut $A + B$. Celui qui gagne paie un prix et garde le lot, donc votre profit est $A + B$ moins ce que vous payez. Si vous perdez, votre profit est nul.

Il y a quatre manches, avec de nouveaux dés à chaque fois. Vous enchérissez par pas de 0.5, de 0 à 20, et en cas d'égalité, c'est votre rival qui l'emporte. Votre rival suit la règle énoncée pour chaque manche, quoi que vous fassiez, et vous connaissez cette règle. Prenez environ 30 secondes pour chaque manche, et essayez de terminer avec le plus d'argent possible.
\begin{itemize}
\item Manche 1. Offres sous pli. L'offre la plus élevée gagne et paie sa propre offre. Votre rival enchérit $B + 1$.
\item Manche 2. Offres sous pli. L'offre la plus élevée gagne et paie l'offre de l'autre joueur. Votre rival enchérit $2B$.
\item Manche 3. Un prix croissant. Le prix monte à partir de 0, et chacun de vous annonce à l'avance le prix le plus élevé auquel il restera. Celui dont la limite est la plus basse se retire, et l'autre gagne à ce prix. La limite de votre rival est $2B$.
\item Manche 4. Comme la manche 1, mais avant d'enchérir, vous pouvez payer 2, que vous gagniez ou non, pour voir le dé $B$.
\end{itemize}

\medskip
\begin{tabular}{@{}l|p{5.5cm}|l|l|l@{}}
Manche & Ce que vous avez décidé & Offre & Résultat & Profit \\ \hline
1 & & & & \\[8pt]\hline
2 & & & & \\[8pt]\hline
3 & & & & \\[8pt]\hline
4 & & & & \\[8pt]\hline
\end{tabular}
\end{question}
\vfill\clearpage

\begin{question}[$\star\star\star\star$ \iqroles{QT, Trader, QR, Risque} \iqfirm{teneur de marché, hedge fund, crypto} \iqlevel{junior, confirmé}]\label{iq:iv2:played-games:4}
\textbf{Cinq paris et un capital de jeu.} Vous partez avec 1\,000. Chaque manche est un pari à une cote donnée, et vous annoncez quelle part de votre capital actuel miser. Vous ne pouvez pas emprunter, et vous pouvez ne rien miser. Une cote annoncée désigne le gain net~: un pari gagnant vous rend également votre mise. Votre objectif est la croissance la plus rapide sur de nombreuses séances de ce type, et non le montant le plus élevé à la fin de celle-ci. Prenez environ 30 secondes par manche.
\begin{itemize}
\item Manche 1. Une pièce qui gagne 60~\% du temps, à mise égale.
\item Manche 2. Un dé équilibré. Un six rapporte 4 fois votre mise. Tout autre résultat fait perdre la mise.
\item Manche 3. Un pari qui gagne 45~\% du temps et rapporte 2 fois votre mise.
\item Manche 4. Deux pièces distinctes, gagnant chacune 60~\% du temps à mise égale. Vous misez le même montant sur chacune.
\item Manche 5. Un pari qui gagne 90~\% du temps et rapporte votre mise. Les 10~\% restants du temps, vous perdez 5 fois votre mise.
\end{itemize}

\medskip
\begin{tabular}{@{}l|l|l|l@{}}
Manche & Fraction misée & Résultat & Capital \\ \hline
1 & & & \\[8pt]\hline
2 & & & \\[8pt]\hline
3 & & & \\[8pt]\hline
4 & & & \\[8pt]\hline
5 & & & \\[8pt]\hline
\end{tabular}
\end{question}
\vfill\clearpage

\begin{question}[$\star\star\star\star$ \iqroles{QT, Trader, QR, Analyste} \iqfirm{teneur de marché, hedge fund} \iqlevel{junior}]\label{iq:iv2:played-games:5}
\textbf{Cotations sur des nombres que vous n'avez pas consultés.} Vous cotez un cours acheteur et un cours vendeur sur cinq nombres, sans calculatrice et pendant environ 30 secondes chacun. Le recruteur connaît chaque nombre. Notez chaque manche ensuite, à la calculatrice~:
\begin{itemize}
\item Vous obtenez toujours $-\log_2(\text{cours vendeur}/\text{cours acheteur})$, de sorte que doubler votre largeur de la fourchette coûte 1.
\item Si le nombre est supérieur à votre cours vendeur, vous obtenez aussi $-4\log_2(\text{nombre}/\text{cours vendeur})$.
\item S'il est inférieur à votre cours acheteur, vous obtenez aussi $-4\log_2(\text{cours acheteur}/\text{nombre})$.
\end{itemize}
Le meilleur score possible est 0. Les nombres sont~: le nombre de nombres premiers inférieurs à 1\,000~; le nombre de chiffres de $100!$~; le nombre de mains de cinq cartes d'un jeu de 52 cartes~; la somme des 100 premiers nombres premiers~; et la somme des chiffres décimaux de $2^{100}$.

\medskip
\begin{tabular}{@{}l|l|l|l|l@{}}
Manche & Cours acheteur & Cours vendeur & Nombre exact & Score \\ \hline
1 & & & & \\[8pt]\hline
2 & & & & \\[8pt]\hline
3 & & & & \\[8pt]\hline
4 & & & & \\[8pt]\hline
5 & & & & \\[8pt]\hline
\end{tabular}
\end{question}
\vfill\clearpage

\begin{question}[$\star\star\star\star$ \iqroles{QT, Trader, QR} \iqfirm{teneur de marché, hedge fund} \iqlevel{junior, confirmé}]\label{iq:iv2:played-games:6}
\textbf{Deux instruments qui partagent un dé.} Trois dés équilibrés à six faces, $A$, $B$ et $C$, seront lancés, dans cet ordre. Deux instruments se règlent sur eux~: $X$ verse $A + B$ et $Y$ verse $A + C$. Chacun a une valeur juste de 7 avant tout lancer et une variance de $35/6$, et leur covariance vaut $35/12$.

Il y a trois manches. Lors de chacune, le recruteur présente des cotations, et vous échangez des lots entiers à ces cotations. Votre responsable juge chaque décision d'après son profit espéré à la valeur juste, diminué de 0.1 fois la variance de la position que vous détenez alors, de sorte qu'un lot qui ajoute de l'avantage peut tout de même être une mauvaise transaction. Seul ce qui est coté peut être négocié. Après la manche 3, tous les dés sont lancés et tout se règle.

\medskip
\begin{tabular}{@{}l|l|l|l@{}}
Manche & Lots de $X$ échangés, à & Lots de $Y$ échangés, à & Position $(X, Y)$ \\ \hline
1 & & & \\[8pt]\hline
2 & & & \\[8pt]\hline
3 & & & \\[8pt]\hline
Règlement-livraison & & & \\[8pt]\hline
\end{tabular}
\end{question}
\vfill\clearpage

\begin{question}[$\star\star\star\star$ \iqroles{QT, Trader, QR, LLE} \iqfirm{teneur de marché, hedge fund, crypto} \iqlevel{junior, confirmé}]\label{iq:iv2:played-games:7}
\textbf{Un contrat, une pièce et un lecteur plus rapide.} Une pièce cachée est de l'un de deux types, chacun étant aussi probable. Une bonne pièce tombe sur face dans 70~\% des cas, et une mauvaise pièce dans 30~\%. Un contrat verse 100 si la pièce est bonne, et 0 si elle est mauvaise.

Il y a six manches. Lors de chacune, vous pouvez soit retirer votre cotation, soit afficher un cours acheteur et un cours vendeur distants de 6, centrés sur la valeur juste que vous annoncez à une décimale près. Si vous affichez une cotation~:
\begin{itemize}
\item un client négocie un lot avec vous, et cela vous garantit 3 par rapport à la valeur juste annoncée, tant que celle-ci est juste~;
\item puis, si un tirage est prévu dans la manche, un lecteur plus rapide le voit avant que vous puissiez modifier votre cotation. Si la nouvelle valeur juste se situe hors de votre cotation, il négocie un lot à votre cotation périmée. Vous clôturez ce lot immédiatement à la nouvelle valeur juste.
\end{itemize}
Le recruteur vous indique à chaque manche si un tirage est prévu. La manche 1 n'en comporte aucun, et les manches 2 à 6 en comportent chacune un. Votre objectif est de gagner le plus d'argent possible.

\medskip
\begin{tabular}{@{}l|l|l|l|l|l@{}}
Manche & Valeur juste & Cours acheteur & Cours vendeur, ou retrait & Tirage & Résultat \\ \hline
\ivscoreline{1}
\ivscoreline{2}
\ivscoreline{3}
\ivscoreline{4}
\ivscoreline{5}
\ivscoreline{6}
\end{tabular}
\end{question}
\vfill\clearpage

\begin{question}[$\star\star\star\star$ \iqroles{QT, Trader} \iqfirm{teneur de marché, crypto} \iqlevel{junior, confirmé}]\label{iq:iv2:played-games:8}
\textbf{Clients et limite de position.} Un contrat vaut exactement 100, et cela ne changera pas. Il y a cinq manches. Lors de chacune, vous annoncez un centre $m$ entier, et vous cotez un cours acheteur de $m - 2$ et un cours vendeur de $m + 2$.

Puis un client arrive avec une valeur privée $v$, chaque entier de 94 à 106 étant aussi probable. Il achète un lot à votre cours vendeur si $v$ dépasse celui-ci, vous vend un lot à votre cours acheteur si $v$ est inférieure à celui-ci, et sinon s'en va. Votre position ne peut jamais dépasser 2 lots, en longue ou en courte~: à la limite, ce côté de votre cotation est simplement désactivé. Après la manche 5, chaque lot que vous détenez encore, en longue ou en courte, vous coûte 4 pour le clôturer. Votre objectif est de maximiser l'espérance de gain.

\medskip
\begin{tabular}{@{}l|l|l|l|l@{}}
Manche & Position avant & Votre $m$ & Client & Trésorerie issue \\ \hline
1 & 0 & & & \\[8pt]\hline
2 & & & & \\[8pt]\hline
3 & & & & \\[8pt]\hline
4 & & & & \\[8pt]\hline
5 & & & & \\[8pt]\hline
\end{tabular}
\end{question}
\vfill\clearpage

\begin{question}[$\star\star\star\star$ \iqroles{QT, Trader, QR} \iqfirm{teneur de marché, hedge fund} \iqlevel{junior, confirmé}]\label{iq:iv2:played-games:9}
\textbf{Un straddle au travers d'une annonce.} Une action vaut 100, sans dérive ni taux d'intérêt. Le jour 1, elle évolue de $+1$ ou de $-1$, chacun de ces mouvements étant équiprobable. Le jour 2, elle évolue de nouveau de $+1$ ou de $-1$, équiprobablement, après une annonce qui la fait d'abord évoluer de $+J$ ou de $-J$, équiprobablement. La taille $J$ vaut 6 ou 10, avec la même probabilité, et personne ne connaît l'une ni l'autre avant l'annonce. Tous ces mouvements sont indépendants.

Un straddle de prix d'exercice 100 verse $|S - 100|$ à la clôture du jour 2. Vous pouvez négocier jusqu'à 2 straddles dans la manche 1, aux cotations affichées, et un nombre entier quelconque d'actions au cours courant dans chaque manche ultérieure. Dans la manche 1, on se trouve avant le jour 1. Dans la manche 2, le jour 1 s'est écoulé. Dans la manche 3, l'annonce a eu lieu, mais le dernier $\pm 1$ n'a pas encore eu lieu.

\medskip
\begin{tabular}{@{}l|p{4cm}|l|l@{}}
Manche & Ce que vous avez fait & Action & Valeur de votre position \\ \hline
1 & & & \\[8pt]\hline
2 & & & \\[8pt]\hline
3 & & & \\[8pt]\hline
Règlement-livraison & & & \\[8pt]\hline
\end{tabular}
\end{question}
\vfill\clearpage

\begin{question}[$\star\star\star\star$ \iqroles{QT, Trader, QR} \iqfirm{teneur de marché, hedge fund, crypto} \iqlevel{junior, confirmé}]\label{iq:iv2:played-games:10}
\textbf{Les cotations des autres joueurs.} Quatre joueurs lancent chacun un dé à dix faces numérotées de 1 à 10, et ne voient que le leur. Un contrat verse la somme des quatre dés. Vous obtenez un 7. Les trois autres suivent des règles fixes, que vous connaissez~:
\begin{itemize}
\item Les joueurs $A$ et $B$ cotent une seule fois, centrés sur leur propre dé plus 16.5, avec un cours acheteur 0.5 en dessous et un cours vendeur 0.5 au-dessus. Chacun négocie un seul lot pendant toute la partie, à cette cotation, et ne consulte jamais celle d'un autre.
\item Le joueur $C$ cote deux fois, centré sur son propre dé plus $n$ plus 16.5. Ici $n$ vaut $-3$, 0 ou $+3$, chacune de ces valeurs étant aussi probable, et tiré à nouveau pour chaque cotation. Le joueur $C$ négocie un lot à chaque cotation.
\end{itemize}
Les cotations apparaissent manche par manche. Vous pouvez négocier ou passer à chacune, et votre seul objectif est de gagner de l'argent. Après la manche 4, les dés sont révélés, et le contrat se règle.

\medskip
\begin{tabular}{@{}l|l|p{5cm}|l@{}}
Manche & Votre valeur juste & Ce que vous avez négocié, avec qui, à quel prix & Position \\ \hline
1 & & & \\[8pt]\hline
2 & & & \\[8pt]\hline
3 & & & \\[8pt]\hline
4 & & & \\[8pt]\hline
\end{tabular}
\end{question}
\vfill\clearpage

\begin{question}[$\star\star\star\star$ \iqroles{QT, Trader, QR, LLE} \iqfirm{teneur de marché, hedge fund, banque} \iqlevel{junior, confirmé}]\label{iq:iv2:played-games:11}
\textbf{Acheter six lots en quatre manches.} Vous devez acheter exactement 6 lots sur quatre manches. Le prix de la manche 1 vaut 100. Dans une manche de prix $P$, l'achat de $k$ lots coûte au total $kP + 0.5k^2$, et cet impact ne se prolonge pas dans la manche suivante.

Au début de chaque manche, on vous communique le prix et une prévision $f$. Le prix de la manche suivante vaut $P + f + \varepsilon$, avec $\varepsilon$ égal à $+1$ ou à $-1$, avec la même probabilité. Les prévisions des manches ultérieures sont inconnues et ont une moyenne nulle. Votre objectif est le coût total espéré le plus faible. La manche 4 doit mener le travail à son terme.

\medskip
\begin{tabular}{@{}l|l|l|l|l@{}}
Manche & Prix & Prévision & Lots achetés & Coût \\ \hline
1 & 100 & & & \\[8pt]\hline
2 & & & & \\[8pt]\hline
3 & & & & \\[8pt]\hline
4 & & & & \\[8pt]\hline
\end{tabular}
\end{question}
\vfill\clearpage

\begin{question}[$\star\star\star\star$ \iqroles{QT, Trader, Risque} \iqfirm{teneur de marché, hedge fund, crypto} \iqlevel{junior}]\label{iq:iv2:played-games:12}
\textbf{Jouer pour la première place.} Vous êtes 3 points derrière un adversaire, qui a cessé de jouer, de sorte que son score est fixé. Il y a quatre manches, chacune jouée à la cote 1 contre 1. Vous misez un nombre entier de points de 0 à 4, et vous gagnez ou perdez autant de points. La probabilité de gagner est de 50~\% dans la manche 1, 55~\% dans la manche 2, 40~\% dans la manche 3 et 60~\% dans la manche 4.

Le seul prix revient à celui qui termine strictement en tête. Une égalité compte comme une défaite. Vos points ne comptent pour rien d'autre. Prenez environ 30 secondes par manche.

\medskip
\begin{tabular}{@{}l|l|l|l|l@{}}
Manche & Écart avant & Votre mise & Résultat & Écart après \\ \hline
1 & $-3$ & & & \\[8pt]\hline
2 & & & & \\[8pt]\hline
3 & & & & \\[8pt]\hline
4 & & & & \\[8pt]\hline
\end{tabular}
\end{question}
